% !TEX program = pdflatex
% !TEX option  = -synctex=1 -interaction=nonstopmode -output-directory=build
% ==========================================================
% BEAMER PRO — PENSAMIENTO MATEMÁTICO (PSICOLOGÍA)
% Tema: Conectivos lógicos, axiomas y análisis de proposiciones
% Autor: Luis Tulio González Alcaino
% Compilación: pdfLaTeX
% ==========================================================

% -------------------------
% MODO
% -------------------------
\newif\ifhandout
\handoutfalse
% \handouttrue

\ifhandout
\documentclass[aspectratio=169,10pt,handout]{beamer}
\else
\documentclass[aspectratio=169,10pt]{beamer}
\fi

% -------------------------
% PAQUETES
% -------------------------
\usepackage[spanish,es-noshorthands]{babel}
\usepackage[T1]{fontenc}
\usepackage[utf8]{inputenc}
\usepackage{lmodern}
\usepackage{amsmath,amssymb,amsfonts}
\usepackage{array,booktabs,multirow}
\usepackage{tabularx}
\usepackage{tikz}
\usetikzlibrary{positioning}
\usepackage{xcolor}
\usepackage{graphicx}
\usepackage{ragged2e}
\usepackage{etoolbox}
\usepackage{enumitem}
\usepackage{pgfplots}
\pgfplotsset{compat=1.18}

% -------------------------
% COLORES — PALETA PSICOLOGÍA
% -------------------------
\definecolor{PsyPrimary}{HTML}{5B4B8A}
\definecolor{PsySecondary}{HTML}{7C6CB3}
\definecolor{PsyAccent}{HTML}{3AA6B9}
\definecolor{PsyLight}{HTML}{F4F1FB}
\definecolor{PsyDark}{HTML}{1F2230}
\definecolor{PsySoft}{HTML}{E7E1F5}
\definecolor{PsyWarn}{HTML}{E67E22}
\definecolor{PsyGood}{HTML}{2E8B57}
\definecolor{PsyBad}{HTML}{C0392B}

% -------------------------
% CONFIGURACIÓN GENERAL
% -------------------------
\setlength{\parskip}{0.35em}

\setbeamercolor{normal text}{fg=PsyDark,bg=white}
\setbeamercolor{structure}{fg=PsyPrimary}
\setbeamercolor{frametitle}{fg=white,bg=PsyPrimary}
\setbeamercolor{title}{fg=white,bg=PsyPrimary}
\setbeamercolor{subtitle}{fg=white}
\setbeamercolor{author}{fg=white}
\setbeamercolor{date}{fg=white}

\setbeamercolor{block title}{fg=white,bg=PsyPrimary}
\setbeamercolor{block body}{fg=PsyDark,bg=PsyLight}

\setbeamercolor{block title alerted}{fg=white,bg=PsyBad}
\setbeamercolor{block body alerted}{fg=PsyDark,bg=red!5}

\setbeamercolor{block title example}{fg=white,bg=PsyGood}
\setbeamercolor{block body example}{fg=PsyDark,bg=green!5}

\setbeamerfont{title}{series=\bfseries,size=\Large}
\setbeamerfont{subtitle}{size=\normalsize}
\setbeamerfont{frametitle}{series=\bfseries,size=\large}

\setbeamertemplate{navigation symbols}{}

% -------------------------
% FOOTLINE MODERNA
% -------------------------
\setbeamertemplate{footline}{
	\leavevmode
	\hbox{
		\begin{beamercolorbox}[wd=.72\paperwidth,ht=2.8ex,dp=1.2ex,leftskip=1em]{author in head/foot}
			\color{PsyDark}\footnotesize \textbf{Pensamiento Matemático} \;--\; Psicología
		\end{beamercolorbox}
		\begin{beamercolorbox}[wd=.28\paperwidth,ht=2.8ex,dp=1.2ex,rightskip=1em plus1fil]{date in head/foot}
			\color{PsyDark}\footnotesize \insertframenumber/\inserttotalframenumber
		\end{beamercolorbox}
	}
}

% -------------------------
% FRAMETITLE
% -------------------------
\setbeamertemplate{frametitle}{
	\nointerlineskip
	\begin{beamercolorbox}[wd=\paperwidth,ht=3.2ex,dp=1.2ex,leftskip=0.8em,rightskip=0.3em]{frametitle}
		\usebeamerfont{frametitle}\insertframetitle
	\end{beamercolorbox}
	\vspace{0.15cm}
}

% -------------------------
% TÍTULO
% -------------------------
\title{Pensamiento Matemático}
\subtitle{Conectivos lógicos, axiomas y análisis del carácter de una proposición}
\author{Profesor Luis Tulio González Alcaino}
\institute{Primer año universitario -- Psicología}
\date{Universidad Santo Tomás}

% -------------------------
% COMANDOS PERSONALIZADOS
% -------------------------
\newcommand{\destacado}[1]{\textcolor{PsyPrimary}{\textbf{#1}}}
\newcommand{\clave}[1]{\textcolor{PsyAccent}{\textbf{#1}}}
\newcommand{\verdadero}{\textcolor{PsyGood}{\textbf{V}}}
\newcommand{\falso}{\textcolor{PsyBad}{\textbf{F}}}

\newenvironment{IdeaClave}[1]
{
	\begin{block}{#1}
		\justifying
	}
	{
	\end{block}
}

\newenvironment{Alerta}[1]
{
	\begin{alertblock}{#1}
		\justifying
	}
	{
	\end{alertblock}
}

\newenvironment{Ejemplo}[1]
{
	\begin{exampleblock}{#1}
		\justifying
	}
	{
	\end{exampleblock}
}

% -------------------------
% PORTADA PERSONALIZADA
% -------------------------
\setbeamertemplate{title page}{
	\begin{tikzpicture}[remember picture,overlay]
		\fill[PsyPrimary] (current page.south west) rectangle (current page.north east);
		\fill[PsySecondary] (current page.south west) rectangle ([xshift=4.2cm]current page.north west);
		\fill[PsyAccent,opacity=0.18] ([xshift=10.8cm,yshift=1.2cm]current page.south west) circle (2.2cm);
		\fill[white,opacity=0.06] ([xshift=12.3cm,yshift=5.8cm]current page.south west) circle (1.5cm);
		\draw[white,opacity=0.10,line width=1.2pt] ([xshift=1.1cm,yshift=1.0cm]current page.south west) -- ([xshift=14.8cm,yshift=1.0cm]current page.south west);
	\end{tikzpicture}
	
	\vspace*{0.8cm}
	
	\begin{columns}[T,totalwidth=\textwidth]
		\begin{column}{0.62\textwidth}
			{\color{white}\usebeamerfont{title}\inserttitle\par}
			\vspace{0.35cm}
			{\color{white}\usebeamerfont{subtitle}\insertsubtitle\par}
			\vspace{0.8cm}
			{\color{white}\normalsize \textbf{\insertauthor}\par}
			\vspace{0.2cm}
			{\color{white}\small \insertinstitute\par}
			\vspace{0.2cm}
			{\color{white}\small \insertdate\par}
		\end{column}
		
		\begin{column}{0.34\textwidth}
			\vspace{0.4cm}
			\begin{tikzpicture}[scale=0.9]
				\draw[white,line width=1.5pt] (0,0) circle (1.0);
				\draw[white,line width=1.2pt] (-0.45,0.25) .. controls (-0.7,0.55) and (-0.25,0.75) .. (0,0.55);
				\draw[white,line width=1.2pt] (0,0.55) .. controls (0.25,0.8) and (0.7,0.55) .. (0.45,0.2);
				\draw[white,line width=1.2pt] (-0.35,-0.05) .. controls (-0.15,0.15) and (0.15,0.15) .. (0.35,-0.05);
				\draw[white,line width=1.2pt] (0,-0.45) -- (0,-0.05);
				\filldraw[white] (-1.1,-1.2) circle (0.07);
				\filldraw[white] (0,-1.55) circle (0.07);
				\filldraw[white] (1.1,-1.15) circle (0.07);
				\draw[white,line width=1pt] (-1.1,-1.2)--(0,-1.55)--(1.1,-1.15);
				\draw[white,line width=1pt] (-0.25,-0.95)--(-1.1,-1.2);
				\draw[white,line width=1pt] (0,-1.0)--(0,-1.55);
				\draw[white,line width=1pt] (0.25,-0.95)--(1.1,-1.15);
			\end{tikzpicture}
			
			\vspace{0.5cm}
			
			\begin{beamercolorbox}[wd=\linewidth,sep=0.4em,rounded=true,center]{block body}
				\color{PsyPrimary}\small
				\textbf{Lógica, análisis y razonamiento}\\
				aplicados al contexto psicológico
			\end{beamercolorbox}
		\end{column}
	\end{columns}
	
	\vfill
}

% -------------------------
% DOCUMENTO
% -------------------------
\begin{document}
	
	% ------------------------------------------------
	\begin{frame}
		\titlepage
	\end{frame}
	
	% ------------------------------------------------
	\begin{frame}{Resultados de aprendizaje}
		\begin{IdeaClave}{Al finalizar la clase, el estudiante será capaz de:}
			\begin{itemize}[leftmargin=*,label=\textcolor{PsyAccent}{\large$\blacktriangleright$}]
				\item Reconocer los \destacado{conectivos lógicos} a través de enunciados gramaticales.
				\item Utilizar los \destacado{axiomas de los conectivos lógicos}.
				\item Emplear esquemas lógicos determinando si una proposición es \destacado{tautología, contradicción o contingencia}.
				\item Integrar \destacado{patrones y relaciones matemáticas} en las actividades realizadas.
				\item Detectar \destacado{patrones, relaciones e inconsistencias} en distintos razonamientos.
			\end{itemize}
		\end{IdeaClave}
	\end{frame}
	
	% ------------------------------------------------
	\begin{frame}{Motivación inicial}
		\begin{Alerta}{Situación para analizar}
			\begin{quote}
				``Si una persona presenta estrés, entonces necesariamente presentará ansiedad.''
			\end{quote}
		\end{Alerta}
		
		\pause
		
		\begin{columns}[T,totalwidth=\textwidth]
			\begin{column}{0.52\textwidth}
				\begin{IdeaClave}{Preguntas orientadoras}
					\begin{itemize}[leftmargin=*]
						\item ¿Es siempre verdadera esa afirmación?
						\item ¿Qué relación lógica aparece?
						\item ¿Puede existir una excepción?
					\end{itemize}
				\end{IdeaClave}
			\end{column}
			\begin{column}{0.44\textwidth}
				\begin{Ejemplo}{Idea central}
					La lógica no decide por sí sola si una afirmación psicológica es empíricamente cierta, pero sí permite \destacado{analizar su estructura de razonamiento}.
				\end{Ejemplo}
			\end{column}
		\end{columns}
	\end{frame}
	
	% ------------------------------------------------
	\begin{frame}{Conectivos lógicos en lenguaje natural}
		\begin{center}
			\renewcommand{\arraystretch}{1.35}
			\begin{tabular}{>{\raggedright\arraybackslash}p{3.2cm} >{\centering\arraybackslash}p{1.7cm} >{\raggedright\arraybackslash}p{5.8cm}}
				\toprule
				\textbf{Conectivo} & \textbf{Símbolo} & \textbf{Expresión habitual} \\
				\midrule
				Negación & $\neg p$ & no, nunca, no es cierto que \\
				Conjunción & $p \land q$ & y, además, junto con \\
				Disyunción & $p \lor q$ & o, o bien \\
				Condicional & $p \to q$ & si... entonces... \\
				Bicondicional & $p \leftrightarrow q$ & si y solo si \\
				\bottomrule
			\end{tabular}
		\end{center}
		
		\pause
		
		\begin{Ejemplo}{Ejemplos contextualizados en Psicología}
			\begin{itemize}[leftmargin=*]
				\item ``El paciente presenta ansiedad \textbf{y} alteraciones del sueño.'' $\Rightarrow p \land q$
				\item ``El paciente \textbf{no} presenta depresión.'' $\Rightarrow \neg p$
				\item ``\textbf{Si} existe estrés crónico, \textbf{entonces} puede disminuir la concentración.'' $\Rightarrow p \to q$
			\end{itemize}
		\end{Ejemplo}
	\end{frame}
	
	% ------------------------------------------------
	\begin{frame}{Traducción de enunciados al lenguaje simbólico}
		\begin{columns}[T,totalwidth=\textwidth]
			\begin{column}{0.48\textwidth}
				\begin{IdeaClave}{Definimos}
					\[
					p:\ \text{``El estudiante duerme bien''}
					\]
					\[
					q:\ \text{``El estudiante mantiene una buena atención''}
					\]
				\end{IdeaClave}
			\end{column}
			\begin{column}{0.48\textwidth}
				\begin{Ejemplo}{Traducción}
					\begin{itemize}[leftmargin=*]
						\item Duerme bien y mantiene buena atención:
						\[
						p \land q
						\]
						\item No duerme bien:
						\[
						\neg p
						\]
						\item Si duerme bien, entonces mantiene buena atención:
						\[
						p \to q
						\]
					\end{itemize}
				\end{Ejemplo}
			\end{column}
		\end{columns}
	\end{frame}
	
	% ------------------------------------------------
	\begin{frame}{Axiomas y relaciones fundamentales}
		\begin{IdeaClave}{Dos esquemas esenciales}
			\[
			p \lor \neg p
			\qquad\qquad
			p \land \neg p
			\]
		\end{IdeaClave}
		
		\pause
		
		\begin{columns}[T,totalwidth=\textwidth]
			\begin{column}{0.48\textwidth}
				\begin{Ejemplo}{Ley del tercero excluido}
					\[
					p \lor \neg p
					\]
					Siempre resulta \destacado{verdadera}. Por eso es una \clave{tautología}.
					
					Ejemplo:
					\begin{quote}
						``Una persona presenta ansiedad o no presenta ansiedad.''
					\end{quote}
				\end{Ejemplo}
			\end{column}
			\begin{column}{0.48\textwidth}
				\begin{Alerta}{Ley de no contradicción}
					\[
					p \land \neg p
					\]
					Siempre resulta \destacado{falsa}. Por eso es una \clave{contradicción}.
					
					Ejemplo:
					\begin{quote}
						``Una persona presenta ansiedad y no presenta ansiedad.''
					\end{quote}
				\end{Alerta}
			\end{column}
		\end{columns}
	\end{frame}
	
	% ------------------------------------------------
	\begin{frame}{Tabla de verdad: conjunción}
		\begin{center}
			\renewcommand{\arraystretch}{1.5}
			\begin{tabular}{c c c}
				\toprule
				$p$ & $q$ & $p \land q$ \\
				\midrule
				\verdadero & \verdadero & \verdadero \\
				\verdadero & \falso     & \falso \\
				\falso     & \verdadero & \falso \\
				\falso     & \falso     & \falso \\
				\bottomrule
			\end{tabular}
		\end{center}
		
		\pause
		
		\begin{IdeaClave}{Interpretación}
			La proposición compuesta $p \land q$ es verdadera \destacado{solo cuando ambas proposiciones son verdaderas}.
		\end{IdeaClave}
	\end{frame}
	
	% ------------------------------------------------
	\begin{frame}{Tabla de verdad: condicional}
		\begin{center}
			\renewcommand{\arraystretch}{1.5}
			\begin{tabular}{c c c}
				\toprule
				$p$ & $q$ & $p \to q$ \\
				\midrule
				\verdadero & \verdadero & \verdadero \\
				\verdadero & \falso     & \falso \\
				\falso     & \verdadero & \verdadero \\
				\falso     & \falso     & \verdadero \\
				\bottomrule
			\end{tabular}
		\end{center}
		
		\pause
		
		\begin{Ejemplo}{Lectura del condicional}
			Si decimos:
			\[
			p:\ \text{``El estudiante estudia''}, \qquad
			q:\ \text{``El estudiante aprueba''}
			\]
			entonces la única situación que hace falsa a $p \to q$ es:
			\begin{quote}
				\textbf{Estudia, pero no aprueba.}
			\end{quote}
		\end{Ejemplo}
	\end{frame}
	
	% ------------------------------------------------
	\begin{frame}{Carácter de una proposición}
		\begin{center}
			\renewcommand{\arraystretch}{1.35}
			\begin{tabular}{>{\raggedright\arraybackslash}p{2.8cm} >{\raggedright\arraybackslash}p{3.3cm} >{\raggedright\arraybackslash}p{5.2cm}}
				\toprule
				\textbf{Tipo} & \textbf{Definición} & \textbf{Ejemplo} \\
				\midrule
				Tautología & Siempre verdadera & $p \lor \neg p$ \\
				Contradicción & Siempre falsa & $p \land \neg p$ \\
				Contingencia & A veces verdadera y a veces falsa & $p \to q$ \\
				\bottomrule
			\end{tabular}
		\end{center}
		
		\pause
		
		\begin{IdeaClave}{Criterio práctico}
			Para clasificar una proposición se analiza su \destacado{tabla de verdad completa}.
			\begin{itemize}[leftmargin=*]
				\item Todas verdaderas $\Rightarrow$ tautología
				\item Todas falsas $\Rightarrow$ contradicción
				\item Mezcla de valores $\Rightarrow$ contingencia
			\end{itemize}
		\end{IdeaClave}
	\end{frame}
	
	% ------------------------------------------------
	\begin{frame}{Ejemplo de clasificación}
		\begin{Ejemplo}{Analicemos la proposición}
			\[
			(p \to q)
			\]
		\end{Ejemplo}
		
		\begin{center}
			\renewcommand{\arraystretch}{1.4}
			\begin{tabular}{c c c}
				\toprule
				$p$ & $q$ & $p \to q$ \\
				\midrule
				\verdadero & \verdadero & \verdadero \\
				\verdadero & \falso     & \falso \\
				\falso     & \verdadero & \verdadero \\
				\falso     & \falso     & \verdadero \\
				\bottomrule
			\end{tabular}
		\end{center}
		
		\pause
		
		\begin{Alerta}{Conclusión}
			Como aparecen valores \destacado{verdaderos y falsos}, la proposición es una \clave{contingencia}.
		\end{Alerta}
	\end{frame}
	
	% ------------------------------------------------
	\begin{frame}{Patrones, relaciones e inconsistencias}
		\begin{columns}[T,totalwidth=\textwidth]
			\begin{column}{0.48\textwidth}
				\begin{IdeaClave}{Patrón lógico}
					Observar regularidades permite anticipar comportamientos de una estructura lógica.
					
					Ejemplo:
					\[
					p \lor \neg p
					\]
					siempre es verdadera.
				\end{IdeaClave}
			\end{column}
			\begin{column}{0.48\textwidth}
				\begin{Alerta}{Inconsistencia}
					Una inconsistencia aparece cuando un razonamiento afirma algo y su negación al mismo tiempo.
					
					Ejemplo:
					\[
					p \land \neg p
					\]
				\end{Alerta}
			\end{column}
		\end{columns}
		
		\pause
		
		\begin{Ejemplo}{Enunciado contextualizado}
			\begin{quote}
				``Todos los pacientes mejoran con terapia''
				
				``Algunos pacientes no mejoran con terapia''
			\end{quote}
			Estas afirmaciones no pueden sostenerse simultáneamente sin revisión, pues generan una \destacado{tensión lógica}.
		\end{Ejemplo}
	\end{frame}
	
	% ------------------------------------------------
	\begin{frame}{Actividad guiada I: reconocer conectivos}
		\begin{IdeaClave}{Indique el conectivo lógico principal de cada enunciado}
			\begin{enumerate}[leftmargin=*,label=\textbf{\arabic*.}]
				\item El paciente está ansioso \textbf{y} presenta insomnio.
				\item El estudiante \textbf{no} manifiesta agotamiento emocional.
				\item \textbf{Si} existe sobrecarga académica, \textbf{entonces} puede disminuir la atención.
				\item El paciente presenta síntomas de estrés \textbf{o} de ansiedad.
				\item Una intervención es efectiva \textbf{si y solo si} mejora el bienestar reportado.
			\end{enumerate}
		\end{IdeaClave}
	\end{frame}
	
	% ------------------------------------------------
	\begin{frame}{Actividad guiada II: simbolización}
		\begin{IdeaClave}{Sea}
			\[
			p:\ \text{``La persona duerme adecuadamente''}
			\qquad
			q:\ \text{``La persona mantiene buena concentración''}
			\]
		\end{IdeaClave}
		
		\begin{Ejemplo}{Escriba simbólicamente}
			\begin{enumerate}[leftmargin=*,label=\textbf{\alph*)}]
				\item La persona duerme adecuadamente y mantiene buena concentración.
				\item La persona no duerme adecuadamente.
				\item Si la persona duerme adecuadamente, entonces mantiene buena concentración.
				\item La persona duerme adecuadamente o mantiene buena concentración.
			\end{enumerate}
		\end{Ejemplo}
	\end{frame}
	
	% ------------------------------------------------
	\begin{frame}{Actividad guiada III: clasificar proposiciones}
		\begin{columns}[T,totalwidth=\textwidth]
			\begin{column}{0.5\textwidth}
				\begin{Ejemplo}{Determine el carácter lógico}
					\begin{enumerate}[leftmargin=*,label=\textbf{\arabic*.}]
						\item $p \lor \neg p$
						\item $p \land \neg p$
						\item $p \to q$
						\item $(p \land q) \to p$
					\end{enumerate}
				\end{Ejemplo}
			\end{column}
			\begin{column}{0.46\textwidth}
				\begin{IdeaClave}{Opciones de clasificación}
					\begin{itemize}[leftmargin=*]
						\item Tautología
						\item Contradicción
						\item Contingencia
					\end{itemize}
				\end{IdeaClave}
				
				\pause
				
				\begin{Alerta}{Sugerencia}
					Antes de responder, identifique el patrón lógico y luego confirme con tabla de verdad.
				\end{Alerta}
			\end{column}
		\end{columns}
	\end{frame}
	
	% ------------------------------------------------
	\begin{frame}{Análisis de un razonamiento cotidiano}
		\begin{Alerta}{Afirmación}
			\begin{quote}
				``Si una persona sonríe, entonces es feliz.''
			\end{quote}
		\end{Alerta}
		
		\pause
		
		\begin{IdeaClave}{Discusión}
			\begin{itemize}[leftmargin=*]
				\item ¿La estructura lógica es condicional? Sí, $p \to q$.
				\item ¿Es necesariamente verdadera en todos los casos? No.
				\item ¿Qué carácter tiene como esquema? \destacado{Contingencia}.
			\end{itemize}
		\end{IdeaClave}
		
		\pause
		
		\begin{Ejemplo}{Reflexión}
			En Psicología, la conducta observable no siempre garantiza un único estado interno; por eso es importante distinguir entre \destacado{apariencia}, \destacado{interpretación} y \destacado{estructura lógica}.
		\end{Ejemplo}
	\end{frame}
	
	% ------------------------------------------------
	\begin{frame}{Síntesis visual de la clase}
		\begin{center}
			\begin{tikzpicture}[
				node distance=1.9cm,
				every node/.style={rounded corners=8pt, align=center, minimum height=1cm, text width=3.0cm, font=\small},
				flecha/.style={->, thick, color=PsyPrimary}
				]
				\node[fill=PsySoft,draw=PsyPrimary] (a) {Conectivos\\lógicos};
				\node[fill=PsySoft,draw=PsyPrimary,right=of a] (b) {Axiomas y\\relaciones básicas};
				\node[fill=PsySoft,draw=PsyPrimary,right=of b] (c) {Tablas de\\verdad};
				\node[fill=PsySoft,draw=PsyPrimary,below=of b] (d) {Tautología,\\contradicción y\\contingencia};
				
				\draw[flecha] (a) -- (b);
				\draw[flecha] (b) -- (c);
				\draw[flecha] (c) -- (d);
				\draw[flecha] (b) -- (d);
			\end{tikzpicture}
		\end{center}
		
		\vspace{0.3cm}
		
		\begin{IdeaClave}{Idea final}
			La lógica permite reconocer estructuras, detectar inconsistencias y fortalecer el razonamiento en contextos académicos y profesionales de la Psicología.
		\end{IdeaClave}
	\end{frame}
	
	% ------------------------------------------------
	\begin{frame}{Evaluación formativa}
		\begin{IdeaClave}{Responde brevemente}
			\begin{enumerate}[leftmargin=*,label=\textbf{\arabic*.}]
				\item ¿Qué conectivo lógico representa la palabra ``y''?
				\item ¿Qué diferencia existe entre tautología y contradicción?
				\item ¿Por qué $p \to q$ suele ser una contingencia?
				\item Escriba un ejemplo simple de inconsistencia lógica.
				\item ¿Cómo ayuda la lógica al análisis de afirmaciones en Psicología?
			\end{enumerate}
		\end{IdeaClave}
	\end{frame}
	
	% ------------------------------------------------
	\begin{frame}{Cierre}
		\begin{Ejemplo}{Conclusiones}
			\begin{itemize}[leftmargin=*,label=\textcolor{PsyAccent}{\large$\checkmark$}]
				\item Los conectivos lógicos permiten traducir lenguaje natural a lenguaje simbólico.
				\item Los axiomas ayudan a reconocer proposiciones siempre verdaderas o siempre falsas.
				\item Las tablas de verdad permiten clasificar proposiciones.
				\item Detectar patrones e inconsistencias fortalece el pensamiento crítico.
			\end{itemize}
		\end{Ejemplo}
		
		\vspace{0.4cm}
		
		\begin{center}
			{\Large \textcolor{PsyPrimary}{\textbf{Pensar con lógica también es comprender mejor la conducta humana.}}}
		\end{center}
	\end{frame}
	
\end{document}